% !TEX root = ../main.tex
% Бүлэг 1

\chapter{Оршил}
\label{Chapter1}
\pagecolor{Beige}

\newcommand{\keyword}[1]{\textbf{#1}}
\newcommand{\tabhead}[1]{\textbf{#1}}
\newcommand{\code}[1]{\texttt{#1}}
\newcommand{\file}[1]{\texttt{\bfseries#1}}
\newcommand{\option}[1]{\texttt{\itshape#1}}

%-------------------------------------------------------------------------------

\section{Судалгааны үндэслэл}

Ясны хугарал нь дэлхий даяар хамгийн түгээмэл тохиолддог гэмтлийн нэг бөгөөд хүн амын эрүүл мэнд, хөдөлмөрийн чадварт ихээхэн нөлөө үзүүлдэг асуудал юм. GBD 2019 судалгааны тайланд дурдсанаар жил бүр дэлхий хэмжээнд 178 сая гаруй шинэ хугарлын тохиолдол бүртгэгддэг байна~\cite{ref01}. Мөн яс-булчингийн тогтолцооны эмгэгүүд нийт 1.71 тэрбум гаруй хүнийг хамарч, хөдөлмөрийн чадвар алдалт болон амьдралын чанар буурах гол шалтгаануудын нэг болж байна~\cite{ref02}.

Монгол Улсад ч мөн адил осол гэмтлийн асуудал нийгмийн эрүүл мэндийн тулгамдсан асуудлын нэг хэвээр байна. 2019 оны байдлаар улсын хэмжээнд 156{,}723 осол гэмтлийн тохиолдол бүртгэгдсэн бөгөөд осол гэмтэл нь нас баралтын тэргүүлэх шалтгаануудын гуравдугаарт эрэмбэлэгдэж байна~\cite{ref03,ref04}. Иймээс ясны хугарлыг богино хугацаанд үнэн зөв оношлох нь эмчилгээний үр дүнг сайжруулах, хүндрэлээс урьдчилан сэргийлэхэд чухал ач холбогдолтой юм.

Сүүлийн жилүүдэд хиймэл оюун ухаан, тэр дундаа конволюцийн нейрон сүлжээ (CNN)-д суурилсан аргууд эмнэлгийн дүрс боловсруулалтын салбарт өргөн хэрэглэгдэх болж, рентген зураг дээрх хугарлыг автоматаар илрүүлэх чиглэлээр өндөр үр дүн үзүүлж байна~\cite{ref08}. Гэсэн хэдий ч хугарлыг илрүүлэх, сегментлэх, хугарлын зэргийг тодорхойлох болон эдгэрэлтийн явцыг үнэлэх зэрэг асуудлуудыг нэгтгэн авч үзсэн судалгаа харьцангуй хязгаарлагдмал хэвээр байна. Иймээс эдгээр чиглэлд ашиглагдаж буй орчин үеийн арга, загваруудыг судлах нь эмнэлгийн дүрс оношилгооны чанар, үр ашгийг дээшлүүлэхэд ач холбогдолтой юм.

%-------------------------------------------------------------------------------

\section{Асуудлын тодорхойлолт}

\subsection*{Уламжлалт оношилгооны хязгаарлалт}

Рентген зураг нь ясны хугарлыг оношлоход хамгийн өргөн хэрэглэгддэг арга боловч нарийн ан цавтай болон бүдэг дүрслэлтэй хугарлуудыг илрүүлэхэд хүндрэлтэй байдаг~\cite{ref10}. Судалгаагаар 3{,}081 батлагдсан хугарлын тохиолдлоос 115 нь анхны үзлэгээр оношлогдоогүй байсан бөгөөд эдгээрийн ихэнхийг нарийн болон рентген зурагт тодорхой ялгарахгүй хугарал эзэлж байжээ~\cite{ref11}. Мөн хугарлын 21.8\% нь анхны рентген зураг дээр илрээгүй болохыг тогтоосон нь уламжлалт оношилгооны хязгаарлалтыг харуулж байна~\cite{ref12}.

\subsection*{Гүн сургалтын аргын өнөөгийн хязгаарлалт}

Гүн сургалтад суурилсан аргууд хугарал илрүүлэх чиглэлээр өндөр үр дүн үзүүлж байгаа боловч судалгааны ихэнх ажил зөвхөн нэг төрлийн даалгаварт төвлөрсөн байдаг~\cite{ref13}. Түүнчлэн чанартай тэмдэглэгээ бүхий өгөгдлийн сан хязгаарлагдмал байгаа нь загваруудын нарийвчлал болон ерөнхийлөх чадварт нөлөөлдөг~\cite{ref14}. Хугарлыг илрүүлэх, сегментлэх болон хугарлын зэргийг тодорхойлох асуудлыг хамтад нь авч үзсэн судалгааны ажил харьцангуй цөөн байна~\cite{ref15}.

\subsection*{Эмнэлгийн үйлчилгээний ачаалал ба хүртээмж}

Ясны хугарлыг оношлох дүрс оношилгооны шинжилгээний тоо нэмэгдэхийн хэрээр радиологич, гэмтлийн эмч нарын өдөр тутмын ажлын ачаалал өсөж байна~\cite{ref16}. Ажлын ачаалал, цаг хугацааны хязгаарлалт болон ядаргаа нь рентген зургийг унших явцад хүний хүчин зүйлээс шалтгаалсан алдаа гарах эрсдэлийг нэмэгдүүлдэг бөгөөд энэ нь ялангуяа нарийн, бүдэг эсвэл шилжилт багатай хугарлын үед илүү тод илэрдэг~\cite{ref07,ref17,ref18}. Радиологийн оношилгоонд алдаа, зөрүү гарах нь тодорхой хэмжээнд зайлшгүй боловч ажлын нөхцөл, ачааллыг сайжруулах болон шийдвэр дэмжих хэрэгсэл ашиглах замаар бууруулах боломжтой гэж үздэг~\cite{ref05}.

Ийм нөхцөлд хиймэл оюунд суурилсан туслах систем нь эмчийн шийдвэрийг орлох бус, харин хугаралтай байж болзошгүй хэсгийг урьдчилан илрүүлэх, сегментлэх, хугарлын зэргийг үнэлэх замаар оношилгооны үйл явцыг дэмжих боломжтой. Дээрх асуудлууд нь ясны хугарлын рентген зурагт гүн сургалтын арга хэрэглэх боломж, хязгаарлалтыг судлах шаардлагатайг харуулж байна.

%-------------------------------------------------------------------------------

\section{Судалгааны зорилго ба зорилт}

\subsection*{Судалгааны зорилго}

Энэхүү судалгааны зорилго нь рентген зурагт суурилан ясны хугарал байгаа эсэхийг ангилах, хугарлын хэсгийг сегментлэх, хугарлын зэргийг үнэлэх болон эдгэрэлтийн хугацааг урьдчилан таамаглах модульчлагдсан хиймэл оюуны pipeline боловсруулахад оршино. Уг зорилгын хүрээнд эхний хоёр туршилтаар нэгдсэн болон тусдаа сургалтын хандлагыг шалгаж, эцсийн туршилтад даалгавар тус бүрийг бие даасан модулиар шийдвэрлэх архитектурыг хэрэгжүүлэн, гүйцэтгэл, давуу тал, хязгаарлалтыг харьцуулан үнэлнэ.

\subsection*{Судалгааны зорилтууд}

Энэхүү зорилгод хүрэхийн тулд дараах зорилтуудыг дэвшүүлэв.

\begin{enumerate}[nosep]
    \item Ясны хугарал илрүүлэх, сегментлэх, хугарлын зэрэг үнэлэх болон эдгэрэлтийн хугацаа таамаглахтай холбоотой онолын үндэс, өмнөх судалгааны ажлуудыг нэгтгэн судалж, энэ ажлын судалгааны цоорхойг тодорхойлох.
    \item MURA, GRAZPEDWRI-DX, FracAtlas-COCO, YOLOBoneFracture болон BoneFractureCVProject зэрэг нийтэд нээлттэй өгөгдлийн сангуудыг судалгааны даалгаварт тохируулан сонгож, ангилал, сегментчлэл, хугарлын зэрэг болон эдгэрэлтийн үнэлгээний зориулалтаар зохион байгуулах.
    \item Рентген зургуудыг сургалтад ашиглахад тохируулан хэмжээг нэг стандарт болгох, нормчлох, өгөгдөл нэмэгдүүлэх, давхардал шалгах, сургалт--баталгаажуулалт--тестийн багцад хуваах урьдчилсан боловсруулалтыг хийх.
    \item Эхний туршилтаар MURA өгөгдөл дээр олон даалгаврыг нэг архитектурт нэгтгэсэн загвар хэрэгжүүлж, даалгавруудыг зэрэг сургах үед гарах давуу болон хүндрэлтэй талыг шинжлэх.
    \item Хоёр дахь туршилтаар ангилал, сегментчлэл, хугарлын зэрэг үнэлгээг тус тусад нь сургах дараалсан сургалтын хандлагыг хэрэгжүүлж, нэгдсэн сургалттай харьцуулах.
    \item Эцсийн туршилтаар ангиллын ConvNeXt-Base модуль, сегментчлэлийн ResNet50 U-Net модуль, хугарлын зэргийн морфологийн шинжид суурилсан модуль болон эдгэрэлтийн таамаглалын модулийг тус тусад нь боловсруулж, нэг pipeline болгон нэгтгэх.
    \item COCO polygon тэмдэглэгээ болон SAM-д суурилсан псевдо маск бэлтгэлийн аргыг ашиглан сегментчлэлийн сургалтын өгөгдлийг бүрдүүлж, хугарлын байршлыг дүрслэх чадварыг үнэлэх.
    \item Сегментчлэлийн маскаас гарган авсан талбай, хэлбэр, периметр зэрэг морфологийн шинжүүдийг ашиглан хугарлын зэргийг үнэлэх боломжийг судлах.
    \item Дүрсний шинж, хугарлын зэрэг болон нас, хүйс зэрэг клиникийн мета-өгөгдлийг нэгтгэн эдгэрэлтийн хугацааг урьдчилан таамаглах туршилтын модулийг боловсруулах.
    \item Боловсруулсан pipeline-ийн үр дүнг ангилал, сегментчлэл, хугарлын зэрэг, эдгэрэлтийн таамаглалын хэмжүүрүүдээр үнэлж, алдааны шалтгаан, тайлбарлах боломж болон эмнэлзүйн хэрэглээнд тавигдах хязгаарлалтыг тодорхойлох.
\end{enumerate}

%-------------------------------------------------------------------------------

\section{Хамрах хүрээ ба хязгаарлалт}

\subsection*{Хамрах хүрээ}

Энэхүү судалгаа нь яс-булчингийн тогтолцооны 2 хэмжээст рентген зурагт суурилсан хугарлын оношилгооны асуудлыг хамарна. Судалгааны хүрээнд хугарлыг ангилах, сегментлэх, хугарлын зэргийг тодорхойлох болон эдгэрэлтийн үнэлгээтэй холбоотой аргуудыг авч үзсэн.

Туршилтад MURA, FracAtlas, GRAZPEDWRI-DX, FracAtlas-COCO, YOLOBoneFracture болон BoneFractureCVProject зэрэг нийтэд нээлттэй өгөгдлийн сангуудыг ашигласан бөгөөд ConvNeXt-Base, ResNet50 U-Net, EfficientNet-B3 зэрэг гүн сургалтын архитектуруудыг туршин судалсан. Судалгааны явцад олон даалгаварт болон бие даасан загваруудын хандлагыг харьцуулан үнэлсэн.

\subsection*{Хязгаарлалт}

\begin{itemize}[nosep]
    \item Судалгаанд зөвхөн 2 хэмжээст рентген зургийг ашигласан бөгөөд КТ, МРТ зэрэг дүрс оношилгооны бусад төрлүүдийг хамруулаагүй.
    \item Ашигласан өгөгдлийн сангуудын тэмдэглэгээ, чанар нь эх өгөгдлийн сангаас шууд хамаарах тул үр дүнд тодорхой хэмжээгээр нөлөөлөх боломжтой.
    \item SAM загвараар үүсгэсэн псевдо шошгууд нь мэргэжлийн эмчийн бүрэн баталгаажуулалтгүй тул сегментчлэлийн чанарт нөлөөлж болзошгүй.
    \item Хугарлын зэргийг тодорхойлоход ашигласан морфологийн шинжүүд нь эмнэлзүйн үнэлгээг бүрэн орлохгүй бөгөөд зөвхөн судалгааны хүрээнд ашиглагдсан.
    \item Эдгэрэлтийн үнэлгээ нь бодит өвчтөний урт хугацааны клиник өгөгдөлд суурилаагүй тул үр дүнг эмнэлзүйн практикт шууд ашиглах боломж хязгаарлагдмал.
    \item Тооцооллын нөөцийн хязгаарлалтын улмаас туршилтуудыг Google Colab болон Kaggle орчинд гүйцэтгэсэн.
    \item Ашигласан өгөгдлийн сангууд нь голчлон гадаад улсын хүн амд суурилсан тул Монгол хүн амд ерөнхийлөн хэрэглэхийн өмнө нэмэлт судалгаа шаардлагатай.
\end{itemize}

%-------------------------------------------------------------------------------

\section{Судалгааны таамаглал}

Энэхүү судалгаанд дараах таамаглалуудыг дэвшүүлэв.

\begin{enumerate}[nosep]
    \item Гүн сургалтад суурилсан ангиллын загвар нь рентген зураг дахь ясны хугарлыг найдвартай ялган таних боломжтой бөгөөд тест нөхцөлд AUC-ROC $0.90$-оос дээш, F1-score $0.85$-аас дээш түвшинд хүрэх боломжтой гэж таамаглав.
    \item Урьдчилан сургасан (pre-trained) архитектур, ялангуяа том хэмжээний дүрсний өгөгдөл дээр сурсан backbone ашиглах нь хугарлын ангиллын гүйцэтгэлийг суурь CNN загвараас сайжруулж, AUC-ROC үзүүлэлтийг $0.95$ орчим буюу түүнээс дээш түвшинд хүргэх боломжтой гэж үзэв.
    \item Хугарлын бүсийг сегментлэх модуль нь хугарлын байршил, хэлбэрийг дүрслэхэд дэмжлэг үзүүлэх бөгөөд тест олонлог дээр Dice coefficient $0.60$ орчим, IoU $0.40$-өөс дээш түвшинд хүрэх боломжтой гэж таамаглав.
    \item Сегментчлэлийн маскаас гарган авсан талбай, периметр, хэлбэрийн харьцаа зэрэг морфологийн шинжүүдийг ашиглан хугарлын зэргийг үнэлэх боломжтой бөгөөд туршилтын нөхцөлд Exact Accuracy $0.85$-аас дээш, adjacent accuracy $0.95$ орчим түвшинд хүрч болзошгүй гэж үзэв.
    \item Хуваалцсан төлөөлөлд суурилсан нэгдсэн олон даалгаварт сургалт нь даалгавруудын хооронд task interference үүсгэж болзошгүй тул бие даасан модульчлагдсан архитектур нь сегментчлэлийн Dice үзүүлэлтийг нэгдсэн загвартай харьцуулахад мэдэгдэхүйц нэмэгдүүлж, ангиллын AUC-ROC-ийг $0.98$ орчим түвшинд хүргэх боломжтой гэж таамаглав.
    \item Дүрсний шинж, хугарлын зэрэг болон клиникийн мета-өгөгдлийг нэгтгэсэн тохиолдолд эдгэрэлтийн хугацааг урьдчилсан түвшинд таамаглах боломжтой бөгөөд регрессийн загварын $R^2$ нь $0.40$-өөс дээш, MAE нь ойролцоогоор 2 долоо хоногийн дотор байх боломжтой гэж үзэв.
    \item Хиймэл оюунд суурилсан pipeline нь оношилгооны эцсийн шийдвэрийг орлохгүй боловч хугарал илрүүлэх магадлал, сегментчлэлийн маск, хугарлын зэрэг болон эдгэрэлтийн урьдчилсан таамаглалыг хамтад нь үзүүлснээр радиологичийн шийдвэр гаргалтыг дэмжих боломжтой гэж таамаглав.
\end{enumerate}
